\section{Discussion and threats to validity}
\label{sec:discussion}

\subsection{What the evidence supports}

The principal result is a layered argument about correction-aware authoritative-state admission. The history and observation model identifies five distinctions; the Alloy model encodes them; ablations expose bounded malformed witnesses when selected encodings are removed; formal--concrete projection binds selected persisted executions to the same relation; and the concrete catalogue checks fail-closed behavior, complete successors, and trace outcomes. The lifecycle makes the dual-value case visible: the machine candidate remains 100 while the human-authorized and committed value is 101.

Each layer licenses a different inference. The paired-history argument explains why a class matters under the declared failure model. Alloy supplies finite sensitivity and separation evidence. Projection checks selected model--database correspondences. The catalogue, stateful tests, and concurrency tests exercise one implementation. The repeated benchmark tests whether agent behavior reaches authority through a restricted surface. Together, the layers connect the information requirements to an executable boundary and its observed behavior.

\subsection{Design implications}

For developers of systems that admit human-corrected AI candidates into persistent records, the first design decision is to preserve the proposal and the authorization as separate immutable objects. The authorization should name the exact candidate and the value permitted to enter the record. This lets a later review distinguish what the machine proposed from what the human approved, even when both lead to the same final value in another history.

Second, select mechanisms by the information they preserve. A stable capability, content-addressed certificate, or database key can represent candidate identity; a version, state digest, or epoch can represent freshness. The paired histories in \cref{tab:distinctions} provide design checks: an implementation needs enough information to separate each relevant safe/unsafe pair, without requiring five physically separate storage objects.

Third, validate the value changes and source evolution together. Atomicity prevents a declared batch from producing a subset of its effects. The total source relation ensures that changed fields acquire their authorized transitions and unchanged fields retain their exact prior sources. A numerically complete snapshot alone is insufficient for this form of accountability.

Finally, make proposal generation and fact admission separate capabilities. Completion and acknowledgment can vary while a trusted host checks the same admission contract. The 720 fixed invalid-tuple calls exercise that contract; the 179 capability-unavailable checks exercise its exposure boundary. These results provide an integration example for the restricted surface. Whether this design improves real reviewers' audit or correction work remains an empirical question for deployments.

\subsection{Formal and construct validity}

SAT and UNSAT describe finite Alloy commands under the stated scopes. The full checks support bounded separation for the encoded catalogue, and conjunct-removal witnesses show sensitivity to selected relational encodings. They do not prove global minimality, unique representation, or unbounded sufficiency. Some high-level classes map to multiple low-level conjuncts; conversely, the source-attribution class is exercised primarily through trace attacks rather than the eleven admission ablations.

The projection contains nine intended instances and twenty mutants, not a proof over all database states. The concrete no-op rule narrows one formal behavior, while SQLite uniqueness excludes one formally representable duplicate-source state. Residual shared-code defects may affect runner, normalizer, model, oracle, or benchmark. Independent projection and database oracles reduce but do not eliminate that risk.

Reverse-trace completeness means that persisted values, sources, certificates, authorizations, and transitions agree under implemented checks. It does not establish evidence truth, candidate accuracy, or sound human judgment. Likewise, Benign Task Completion and Unauthorized Authoritative Mutation measure different constructs; utility failure is not recoded as an authority violation, and runtime failure is reported separately.

\subsection{Concrete and performance validity}

The production-facing catalogue is finite, and concurrency evidence covers the declared SQLite schedules rather than every interleaving or database engine. The trusted implementation includes canonical serialization, hashing, the admission service, the database adapter, configured transaction semantics, and the principal-policy source. Porting the contract requires a new adapter and conformance evidence for the target engine.

The historical admission comparator writes an empty source map and remains only a lower bound. The persistence-equivalent arm reproduces the full relational post-state through a generic delta materializer. Its latency gap combines excluded validation/planning work with implementation-path differences and is not a causal estimate for one check. The optimized trace and baseline use the same grid in separate executions; query-count reduction follows from the access shape, while latency ratios may still reflect caching, scheduling, checkpointing, and runtime noise.

Principal-policy evidence is also narrower than general revocation safety. The adapter rechecks the prototype allow-list inside the transaction before its compare-and-swap. Revocation already visible at that point rejects without authority writes, but revocation after the check does not cancel the in-flight transaction. The prototype therefore demonstrates commit-entry authorization revalidation, not linearizable or distributed revocation.

The lifecycle uses synthetic/public input and a scripted principal. Usability, reviewer error, authentication workflows, other databases, operational workloads, deletion, redaction, retention expiry, and tombstones require separate study. Manual entry remains a compatibility path with a derived certificate and is outside the AI-derived candidate claim.

\subsection{Repeated-agent validity}

The Final benchmark covers three qualified configurations from two provider families, three prompt variants, fourteen scenarios, and ten repetitions per cell. It is broader than the earlier one-off probe but remains a finite convenience population. Configuration identifiers name requested model endpoints and reasoning settings, not immutable model weights. Provider revisions, service load, decoding behavior, and network conditions may change.

Ninety-three of 1,260 planned executions ended in runtime failure, unevenly distributed across configurations (\cref{tab:runtime-endpoint}). Pilot qualification is a pre-Final selection gate, not a predictor of final runtime stability: D1 had 0/28 pilot failures and 76/420 Final failures, while G1 and G2 remained at 4/420 and 13/420. All configurations are retained because the locked protocol forbids post-hoc removal after Final start. Endpoint evaluability is not identical to runtime success: 64 runtime-failure rows retained complete authority verdicts and three retained behavior verdicts. The study was not designed or powered for provider ranking; descriptive differences must not be read as stable model traits. The post-hoc semantic audit corrects the lexical proxy but is still deterministic single-rubric coding rather than independent human annotation. No unavailable-tool attempt was observed, so that measure does not itself demonstrate behavioral heterogeneity.

The authority endpoint covers 720 fixed invalid-tuple admission calls and 179 capability-unavailable branches. Its challenge parameters were not generated by the models. The benchmark therefore does not evaluate model-generated attack search, compromised hosts or identity providers, administrator access, or model access to fact admission. Zero observed violations describe the executed checks, not a production failure rate or general prompt-injection security; the optional binomial reference calculation is documented in \cref{app:zero-event}.

\subsection{Lifecycle, reproducibility, and novelty validity}

The artifact retains immutable source, locked configurations, raw attempts, normalized inputs, generated tables, and hash-linked manifests. The public deposit and verification procedure in \cref{app:artifact} make frozen-record inspection possible without repeating hosted-model calls. Independent reproduction in another environment remains a separate validation step.

The related-work comparison is bounded by the public disclosures available at its stated cutoff. The contribution concerns correction-aware field admission and its declared failure-distinguishability model; the transaction, authorization, freshness, and provenance mechanisms it builds on are established foundations.
